% !TeX spellcheck = es_ES
\documentclass[letterpaper,11pt]{article}
\usepackage[utf8]{inputenc}
\usepackage[T1]{fontenc}
\usepackage[spanish]{babel}
\decimalpoint % separador decimal con punto, como en las tablas de la fuente original
\usepackage{lmodern,microtype}
\usepackage{amsmath,amssymb,amsfonts}
\usepackage{geometry}
\geometry{margin=2.54cm}
\usepackage{setspace}
\setstretch{1.15}
\usepackage{enumitem}
\usepackage{booktabs}
\usepackage{color}
\usepackage{hyperref}
\definecolor{ntr}{rgb}{0.8,0,0}
\hypersetup{colorlinks=true,linkcolor=ntr,citecolor=ntr,urlcolor=black}

% Portada (Luis Cabral): titulo \LARGE sans-serif alineado a la izquierda, sin numero de pagina
\makeatletter
    \renewcommand{\@maketitle}{%
    \newpage\parindent0in\null\vskip2em%
    {\LARGE\textbf{\textsf{\@title}}\par\vskip2em}%
    \@author\par\vskip2em%
    \@date%
    \par
    \thispagestyle{empty}\setcounter{page}{1}%
    }%
\makeatother

\DeclareMathOperator{\Var}{Var}
\DeclareMathOperator{\Cov}{Cov}
\DeclareMathOperator{\Corr}{Corr}
\DeclareMathOperator{\E}{\mathbb{E}}
\DeclareMathOperator{\N}{Normal}
\DeclareMathOperator{\LP}{\mathbb{L}}
\newcommand\independent{\protect\mathpalette{\protect\independenT}{\perp}}
\def\independenT#1#2{\mathrel{\rlap{$#1#2$}\mkern2mu{#1#2}}}
% Teoria asintotica
\newcommand{\SR}{\mathbb{R}}
\DeclareMathOperator{\plim}{plim}
\DeclareMathOperator{\Avar}{AVar}
\newcommand{\pto}{\stackrel{p}{\longrightarrow}}
\newcommand{\dto}{\stackrel{d}{\longrightarrow}}
\newcommand{\adistr}{\stackrel{a}{\sim}}
\newcommand{\epsi}{\varepsilon}


\title{Análisis Econométrico \\ Ayudantía 6}
\author{Magíster en Economía, Universidad Alberto Hurtado}
\date{}

% Funciones de R sugeridas: entre corchetes, alineadas a la derecha; si no caben en la
% linea pasan enteras a la siguiente.
\newcommand{\fn}[1]{{\unskip\nobreak\hfill\penalty50\hskip1em\null\nobreak\hfill
    \mbox{[#1]}\parfillskip=0pt\finalhyphendemerits=0\par}}

\begin{document}
\maketitle

\noindent Temas: unidades de medida, modelos cuadráticos, interacciones entre una variable
binaria y una continua, y efectos marginales.

\medskip
\noindent Los ejercicios 1 y 2 se resuelven en la pizarra. Los ejercicios 3 y 4 son
una demostración en R: el ayudante corre el código en pantalla, así que no hace falta traer
computador. El ejercicio 3 usa los mismos datos de la Ayudantía 5.

\section*{Si quiere replicarlo en su computador}

\noindent El script completo, \texttt{TA6.R}, y los datos están en el repositorio del curso,
en \href{https://github.com/ramirode/econometria_mae/tree/main/TA/TA6}{\texttt{TA/TA6}}. Los
pasos son los mismos de la Ayudantía 5: cree una carpeta \texttt{TA6} con las subcarpetas
\texttt{code} (para \texttt{TA6.R}) e \texttt{input} (para \texttt{ch8\_cps.dta}), y cambie la
línea \texttt{dir\_proyecto <- "..."} por la ruta a su carpeta \texttt{TA6}. En Posit Cloud la
carpeta ya existe: \texttt{/cloud/project/TA/TA6}. Usa los mismos paquetes de la Ayudantía 5.

\section*{Ejercicios}

\begin{enumerate}[itemsep=0.8em]

%% ---------------------------------------------------------------- 1
\item \emph{Unidades de medida.} Se estima por MCO el modelo
    $y_i=\beta_0+\beta_1x_i+u_i$, con estimaciones $\hat\beta_0$ y $\hat\beta_1$, residuos
    $\hat u_i$ y coeficiente de determinación $R^2$.
    \begin{enumerate}
        \item Se cambian las unidades de la variable dependiente: $y_i^*=c\,y_i$, con $c>0$
              (por ejemplo, salarios en miles de pesos en vez de pesos, $c=1/1000$). Encuentre
              las estimaciones MCO de la regresión de $y^*$ en $x$ en función de
              $\hat\beta_0$ y $\hat\beta_1$. ¿Qué pasa con los residuos y con el $R^2$?
        \item Se cambian las unidades de la variable explicativa: $x_i^*=a\,x_i$, con
              $a\neq 0$. ¿Qué pasa con $\hat\beta_0$, $\hat\beta_1$, los residuos y el $R^2$?
        \item Suponga ahora que el modelo es $\log(y_i)=\beta_0+\beta_1x_i+u_i$. ¿Qué pasa con
              las estimaciones si $y$ se mide en otras unidades, $y_i^*=c\,y_i$? ¿Por qué
              cambia sólo la constante?
    \end{enumerate}

%% ---------------------------------------------------------------- 2
\item \emph{Modelo cuadrático.} Considere
    $\E(y|x)=\beta_0+\beta_1x+\beta_2x^2$.
    \begin{enumerate}
        \item Encuentre el efecto marginal de $x$ sobre $\E(y|x)$. ¿Depende de $x$?
        \item Suponga $\beta_1>0$ y $\beta_2<0$. Encuentre el valor de $x$ en que $\E(y|x)$
              alcanza su máximo.
        \item Si la variable dependiente es $\log(y)$, ¿cómo se interpreta el efecto
              marginal de (a)?
    \end{enumerate}

%% ---------------------------------------------------------------- 3
\item \emph{Interacción entre sexo y educación (R).} Considere el modelo
    \begin{align*}
    lahe_i = \beta_0 + \beta_1\, female_i + \beta_2\, yrseduc_i
             + \beta_3\, female_i\times yrseduc_i + u_i,
    \end{align*}
    donde $lahe$ es el logaritmo del salario por hora, $female$ es igual a 1 si la persona es
    mujer y $yrseduc$ son los años de educación.
    \begin{enumerate}
        \item Estime el modelo. Escriba $\E(lahe|female,yrseduc)$ para hombres y para mujeres
              en función de los $\beta$'s. ¿Cuál es el retorno de un año de educación para
              cada sexo?
              \fn{\texttt{lm}, \texttt{coef}}
        \item Estime la regresión de $lahe$ en $yrseduc$ por separado para hombres y para
              mujeres. Compare las constantes y las pendientes con las del modelo con
              interacción.
              \fn{\texttt{filter}, \texttt{modelsummary}}
        \item Use el modelo con interacción para calcular la brecha de salarios entre mujeres
              y hombres con 12 años de educación (secundaria completa) y con 16 años
              (universitaria completa). Exprésela en porcentaje de forma exacta. ¿La brecha
              aumenta o disminuye con la educación?
              \fn{\texttt{exp}}
        \item Agregue al modelo con interacción la edad ($age$) y la edad al cuadrado dividida por 100
              ($age^2/100$). Calcule el efecto marginal de la edad a los 30 y a los 50 años,
              y la edad en que el log salario predicho alcanza su máximo.
              \fn{\texttt{mutate}, \texttt{lm}, \texttt{coef}}
    \end{enumerate}

%% ---------------------------------------------------------------- 4
\item \emph{Errores estándar con heteroscedasticidad (R).} Simule $R=2000$ muestras de
    tamaño $N=200$ del modelo
    \begin{align*}
    y_i=1+2x_i+u_i,\qquad x_i\sim\text{Uniforme}(0,2),\qquad u_i=x_i^2\,e_i,\qquad
    e_i\sim\N(0,1),
    \end{align*}
    con $e_i$ independiente de $x_i$. En cada muestra, estime la regresión de $y$ en $x$ y
    guarde $\hat\beta_1$, su error estándar clásico (que supone homoscedasticidad) y su error
    estándar robusto.
    \begin{enumerate}
        \item ¿Es heteroscedástico el error? Calcule $\E(u_i^2|x_i)$.
        \item Compare la desviación estándar de $\hat\beta_1$ entre las $R$ muestras con el
              promedio de cada error estándar. ¿Cuál se acerca más? ¿Cuál subestima la
              variabilidad de $\hat\beta_1$?
              \fn{\texttt{set.seed}, \texttt{runif}, \texttt{rnorm}, \texttt{vcov},
              \texttt{vcovHC}}
    \end{enumerate}
\end{enumerate}

\end{document}
